\section{Method}
\label{sec:methods}

Our method has two parts that share the same detector scores: an audit design that decides which training labels to acquire, and a score-fixed replay that measures how the negative calibration pool sets the patient alert threshold. Algorithm~\ref{alg:audit} summarizes both.

\subsection{Problem setup}
\label{sec:setting}

Each patient contributes one CT scan together with a fixed \emph{roster} of candidate locations produced by existing detection systems. For a candidate $i$ we denote by $X_i$ the image features extracted around its location. Every scan was annotated by four LIDC readers, whom we randomly assign to the roles A, B, C and D within each scan. Readers A and B provide the \emph{weak reads}: for every candidate we record whether each of them marked a nodule at that location, reviewed the scan without marking it, or was unavailable. Readers C and D define the \emph{reference}. A \emph{target} $j$ is a nodule contoured by both C and D at corresponding locations, and its centre and diameter are the means of their two contours. A candidate receives the reference label $Y_i=1$ if it lies within one target radius, but no less than 3\,mm, of a target, and $Y_i=0$ otherwise. An \emph{omission} is a target that reader A did not mark. Omission status is determined before any candidate is scored, so a target can be an omission even when no candidate covers it.

\subsection{Paired-tail audit design}
\label{sec:audit}

An \emph{audit} reveals the label $Y_i$ of a single training candidate. We assign every candidate an inclusion probability $\pi_i$ and audit it through an independent draw $R_i\sim\mathrm{Bernoulli}(\pi_i)$; known inclusion probabilities allow unbiased corrections later on~\cite{horvitz1952}. The design combines alert leverage, reader disagreement and a budget constraint.

\paragraph{Alert leverage.}
A candidate can only affect a patient alert if it can become the highest score in its scan or exceed the threshold. We therefore freeze a pilot detector, trained on development data with the marks of reader A, which provides a logit score $s_i$, a low-alert threshold $t_0$ on the same logit scale and the rank $r_i\in\{1,2,\dots\}$ of each candidate within its scan, with $r_i=1$ for the highest-scoring one. With a rank scale $\tau>0$, a score offset $\delta$ and a width $\omega>0$, the leverage of candidate $i$ is
\begin{equation}
 h_i=\max\!\left\{e^{-(r_i-1)/\tau},\ \sigma\!\left(\tfrac{s_i-t_0+\delta}{\omega}\right)\right\}\in(0,1],
 \label{eq:leverage}
\end{equation}
where $\sigma$ denotes the logistic function. The first term equals one for the top candidate of each scan and decays with rank; the second approaches one as the score nears the threshold, rising from $\delta$ logits below $t_0$.

\paragraph{Reader disagreement.}
Let $m_i\in\{0,1,2\}$ count the weak marks on candidate $i$, and let $p_i\in[0,1]$ be the out-of-fold prediction of an image model trained to predict whether A or B marked the location. We define a working probability that the candidate is a target and its Bernoulli variance,
\begin{equation}
 u_i=\rho_{m_i}+\eta_{m_i}\,p_i,\qquad v_i=\max\{u_i(1-u_i),\,v_{\min}\},
 \label{eq:working}
\end{equation}
where the intercepts $\rho_m$ and slopes $\eta_m$ depend on the number of marks and the floor $v_{\min}>0$ keeps clear negatives in the design. The coefficients make two concordant marks nearly certain, so that $v_i$ is small, and a single mark highly uncertain, so that $v_i$ approaches its maximum of $1/4$.

\paragraph{Allocation.}
Given an expected budget $\mathcal B=\sum_i\pi_i$, we spend audits where both leverage and uncertainty are high by minimizing $\sum_i \bar h_i^2 v_i\,(1/\pi_i-1)$ with $\bar h_i=\max(h_i,h_{\min})$. Each term is a surrogate for the variance of an inverse-probability estimate of a label with uncertainty $v_i$, weighted by its squared leverage.

\begin{proposition}[Square-root allocation]
\label{prop:alloc}
Among $\pi\in[\pi_{\min},\pi_{\max}]^N$ with $\sum_i\pi_i=\mathcal B$, the objective is minimized by
\begin{equation}
 \pi_i=\operatorname{clip}\!\left(c\,\bar h_i\sqrt{v_i},\ \pi_{\min},\ \pi_{\max}\right),
 \label{eq:allocation}
\end{equation}
for the unique $c>0$ that meets the budget.
\end{proposition}
\begin{proof}
The objective is strictly convex in $\pi$, since $1/\pi_i$ is. At an interior coordinate, the Karush--Kuhn--Tucker condition $\bar h_i^2v_i/\pi_i^2=\lambda$ gives $\pi_i=\bar h_i\sqrt{v_i}/\sqrt\lambda$; coordinates at a bound are clipped, and $c=\lambda^{-1/2}$ is unique because the clipped sum increases monotonically in $c$.
\end{proof}

\noindent We find $c$ by bisection. The budget factor $b\in\{0.5,1\}$ selects the half or base budget, and the floor $\pi_{\min}$, which scales with $b$, gives every candidate a chance of review.

\paragraph{Budget-matched controls.}
\emph{Image-tail} sampling applies Eq.~\eqref{eq:allocation} with a constant $v_i$ and thus ranks candidates without using the marks. \emph{Mark stratification}, the conventional design, sets $\pi_i=(0.02,0.35,0.20)\,b$ for $m_i=0,1,2$, and its expected count defines the budget $\mathcal B$ shared by all policies. \emph{Uniform} sampling sets $\pi_i=\mathcal B/N$ for $N$ candidates. All policies share the same random numbers within each of six matched draws.

\paragraph{Reach.}
Since several candidates may cover the same target, we count distinct targets. Denoting by $\mathcal I_j$ the candidates on target $j$, the probability that at least one of them is audited is
\begin{equation}
 a_j=1-\textstyle\prod_{i\in\mathcal I_j}(1-\pi_i),
 \label{eq:reach}
\end{equation}
with $a_j=0$ when $\mathcal I_j$ is empty. Summing $a_j$ over the omissions gives the expected number of omissions reached.

\paragraph{Implementation constants.}
All constants were fixed before any reference label was revealed. We use $\tau=6$, so that the rank term halves about every four ranks, $\delta=1$ and $\omega=0.7$ in Eq.~\eqref{eq:leverage}; $(\rho_m,\eta_m)=(0,0.35)$, $(0.55,0.35)$ and $(0.92,0.06)$ for $m=0,1,2$ and $v_{\min}=0.006$ in Eq.~\eqref{eq:working}; and $h_{\min}=0.002$, $\pi_{\min}=0.005b$ and $\pi_{\max}=0.95$ in Eq.~\eqref{eq:allocation}. For instance, the fourth-ranked candidate of a scan scoring one logit below $t_0$ obtains $h_i=\max\{0.61,0.50\}=0.61$, whereas the 200th candidate has a leverage close to zero.

\begin{algorithm}[t]
\caption{Paired-tail audit and score-fixed reference replay}
\label{alg:audit}
\small
\begin{algorithmic}[1]
\Require training roster with pilot scores $s_i$, ranks $r_i$, mark counts $m_i$, image predictions $p_i$; budget $\mathcal B$
\State $h_i\gets$ Eq.~\eqref{eq:leverage}; \ $u_i,v_i\gets$ Eq.~\eqref{eq:working}
\State bisect $c$ so that $\sum_i\pi_i=\mathcal B$, with $\pi_i$ from Eq.~\eqref{eq:allocation}
\State draw $R_i\sim\mathrm{Bernoulli}(\pi_i)$; reveal $Y_i$ where $R_i=1$
\State fit $q_i$ out of fold; $z_i\gets q_i+R_i(Y_i-q_i)/\pi_i$
\State fit image-only head $g$ with logistic loss on $z_i$
\For{negative pool $\mathcal N\in\{P,U\}$} \Comment{same scores}
\State $S_j\gets\max_{i\in j} g(X_i)$ for every $j\in\mathcal N$
\State $k\gets\lceil(|\mathcal N|+1)(1-\alpha)\rceil$; \ $T_{\mathcal N}\gets S_{(k)}$
\State count targets and alerts with $g(X_i)>T_{\mathcal N}$
\EndFor
\State compare $T_U$ with 8,192 random removals of $|P|-|U|$ patients
\end{algorithmic}
\end{algorithm}

\subsection{Learning from audited labels}
\label{sec:training}

All deployed detectors operate on image features alone; weak marks and audits are used only for training. A linear logistic head $g$ maps the features $X_i$ to a logit. Because most training candidates are never audited, we fit an auxiliary prediction $q_i$ of $Y_i$ on the other patient folds, using the image embedding, the weak marks, their agreement and their availability, and train on the corrected target
\begin{equation}
 z_i=q_i+\frac{R_i}{\pi_i}\,(Y_i-q_i).
 \label{eq:augmentation}
\end{equation}
Unaudited candidates retain the auxiliary prediction, whereas audited ones add the residual scaled by $1/\pi_i$. We refer to this estimator as \emph{discriminative augmentation}; its key property is the following.

\begin{proposition}[Unbiased training loss]
\label{prop:unbiased}
Let $\ell(y,t)=-y\log\sigma(t)-(1-y)\log(1-\sigma(t))$ be the logistic loss. If $R_i\sim\mathrm{Bernoulli}(\pi_i)$ with $\pi_i>0$, independently of $q_i$, then for any fixed head $g$, $\mathbb E_R\big[\sum_i\ell(z_i,g(X_i))\big]=\sum_i\ell(Y_i,g(X_i))$.
\end{proposition}
\begin{proof}
$\mathbb E[z_i]=q_i+(\pi_i/\pi_i)(Y_i-q_i)=Y_i$, and $\ell$ is affine in its first argument.
\end{proof}

\noindent Independence holds because $q_i$ is fitted on the other patient folds. The expected training loss therefore equals the loss under full reference labels, at the cost of the variance $(1-\pi_i)(Y_i-q_i)^2/\pi_i$, which is large for candidates with a small inclusion probability and an inaccurate auxiliary.

\subsection{Patient-level alert threshold}
\label{sec:cutoff}

The score of patient $j$ is the maximum logit over its candidates, $S_j=\max_{i\in j}g(X_i)$. Given $n$ negative calibration patients with ordered scores $S_{(1)}\le\dots\le S_{(n)}$, the threshold for a target alert rate $\alpha$ is
\begin{equation}
 T_\alpha=S_{(k)},\qquad k=\lceil(n+1)(1-\alpha)\rceil ,
 \label{eq:cutoff}
\end{equation}
and a patient alerts if any of its candidates scores strictly above $T_\alpha$. If the negative maxima are exchangeable with those of future negative patients, the alert probability is $(n+1-k)/(n+1)\le\alpha$~\cite{angelopoulos2023conformal}. In our setting, $n=\PCal$ and $\alpha=0.10$ give $k=\lceil 61.2\rceil=\PRank$, so the threshold is the sixth-largest negative maximum and the alert rate is at most $6/68=8.8\%$. Below, $\alpha=0.10$ is fixed and we write $T_P$ and $T_U$, or simply $T$, for thresholds computed on a negative pool. A target counts as \emph{recovered} if any candidate on it exceeds the threshold.

\subsection{Score-fixed replay of the negative pool}
\label{sec:replay}

The guarantee above holds only for the population that the calibration negatives represent. Let $P$ denote the patients negative under the C/D reference, and $U\subseteq P$ those that are also negative under the official LUNA16 reference, which accepts nodules marked by at least three of the four readers~\cite{setio2017}. The excluded set $E=P\setminus U$ thus contains scans with a nodule that C and D did not both contour. Their maxima can be large, and since $P$ is a mixture of $U$ and $E$, they can raise the upper order statistic. We recompute Eq.~\eqref{eq:cutoff} on $U$ while keeping every candidate score and every original test target, and count alerts on the same $U$-negative test patients under both thresholds, so that any change in recovery is attributable to the threshold alone.

Removing patients also changes $n$ and $k$; with $n=\UCal$, we obtain $k=\lceil 48.6\rceil=\URank$ and an alert rate of at most $5/54=9.3\%$. To separate the effect of \emph{which} patients are removed from that of a smaller pool, we draw 8,192 random subsets of the same size, stratified by calibration subset, and locate each resulting threshold on the original rank scale. We further rebuild the negative pool for each of the six reader pairs AB through CD, keeping the original targets and a test-negative set common to all pairs.
